\subsection{Грузики на веревке}
%%%%%%%%%%%%%%%%%%%%%%%%%%%%%%%%%%%%%%%%%%%%%%%%%%%%
%%%%%%%%%%%%%%%%%%%%%%%%%%%%%%%%%%%%%%%%%%%%%%%%%%%% 
\z[{\begin{wrapfigure}[7]{r}{0.3\textwidth}%
		\vspace{-10pt}%
		\begin{flushright}%
 			\includegraphics[width=0.3\textwidth]{7.4.1}%
		\end{flushright}%
	\end{wrapfigure}%
	}]%
	{На горизонтальной поверхности стола находится платформа с укреплённым на ней штативом. К штативу привязан на нити длиной $l$ небольшой по сравнению с длиной нити шар. Масса платформы со штативом $2m$, масса шара $3m$. Шар отклоняют и удерживают неподвижно так, что нить составляет угол $\gamma$ ($\cos \gamma = \frac{1}{3}$) с вертикалью, а платформа прижата к упору (см. рисунок). Затем шар отпускают.
	\lpu Найдите скорость шара в момент отрыва платформы от упора. \lpuend
	\lpu Найдите максимальный угол отклонения нити от вертикали направо в процессе движения системы после отрыва от упора. \lpuend
	}
%
{\puan  $v = \sqrt{\frac{4}{3}gl};$~\puan $\alpha = \arccos \frac{11}{15}$}
%%%%%%%%%%%%%%%%%%%%%%%%%%%%%%%%%%%%%%%%%%%%%%%%%%%%	
\z[{\begin{wrapfigure}[6]{r}{0.2\textwidth}%
		\vspace{-25pt}%
		\begin{flushright}%
 			\includegraphics[width=0.2\textwidth]{7.4.2}%
		\end{flushright}%
	\end{wrapfigure}%
	}]%
	{Два шара -- стальной массы $m$ и свинцовый массы $\frac{m}{4}$ -- подвешены в одной точке на нитях длины $L$ каждая (см. рисунок). Свинцовый шар отклоняют так, что нить образует угол $\alpha$ с вертикалью, и отпускают. После соударения шар отклоняется на угол $\beta$. Удар центральный. Определите энергию $\Delta E$, перешедшую в тепло.}
%
{$\Delta E = \frac{1}{8}mgL\left(3\sin^2\frac{\alpha}{2} - 5\sin^2\frac{\beta}{2} + 2\sin\frac{\alpha}{2}\sin\frac{\beta}{2}\right)$}
%%%%%%%%%%%%%%%%%%%%%%%%%%%%%%%%%%%%%%%%%%%%%%%%%%%% 
 \z{Маленький шарик массой $m$ висит неподвижно на невесомой нерастяжимой нити длиной $l$. Шарику толчком сообщают такую горизонтальную скорость, что при последующем движении шарик поднимается над начальной точкой на высоту, меньшую $l$, а минимальная сила натяжения нити равна $\frac{mg}{3}$. На какой высоте находился шарик в момент, когда сила натяжения нити равнялась $mg$?}
%
{$h = \frac{l}{3}$}
%%%%%%%%%%%%%%%%%%%%%%%%%%%%%%%%%%%%%%%%%%%%%%%%%%%%
\z[{\begin{wrapfigure}[5]{r}{0.3\textwidth}%
		\vspace{-10pt}%
		\begin{flushright}%
 			\includegraphics[width=0.3\textwidth]{7.4.4}%
		\end{flushright}%
	\end{wrapfigure}%
	}]%
	{На гладкой наклонной плоскости с углом наклона $\alpha$ к горизонту в точке $O$ прикреплена нить длиной $l$. К другому концу нити привязан небольшой шарик (см. рисунок). В начальный момент шарик находится в положении равновесия в точке $A$. Какую минимальную скорость надо сообщить шарику в точке $A$ вдоль наклонной плоскости в горизонтальном направлении, чтобы шарик совершил полный оборот, двигаясь по окружности?}
%
{$v = \sqrt{5gl\sin\alpha}$}
%%%%%%%%%%%%%%%%%%%%%%%%%%%%%%%%%%%%%%%%%%%%%%%%%%%%	
\z[{\begin{wrapfigure}[6]{r}{0.5\textwidth}%
		\vspace{-10pt}%
		\begin{flushright}%
 			\includegraphics[width=0.5\textwidth]{7.4.5}%
		\end{flushright}%
	\end{wrapfigure}%
	}]%
	{С высоты $\frac{3R}{2}$ соскальзывает без начальной скорости небольшой шарик, двигаясь без трения по жёлобу, расположенному в вертикальной плоскости (см. рисунок). Горизонтальный участок жёлоба плавно переходит в полуокружность радиуса $R$.
	\lpu Какой максимальной высоты $H$ достигнет шарик после отрыва от жёлоба? \lpuend
	\lpu Под каким углом $\beta$ к горизонту упадёт шарик на горизонтальный участок жёлоба? \lpuend
}
%
{\puan  $H = \frac{40}{27}R;$~\puan $\beta = \arccos\left(\frac{1}{9}\right)$}